\documentclass[a4paper,11pt,reqno]{amsart}

\usepackage[T1]{fontenc}
\usepackage[expansion=false]{microtype}
\usepackage{amsmath,amssymb,amsthm,mathtools}
\usepackage{xcolor}
\usepackage{enumitem}
\usepackage[a4paper,left=22mm,right=22mm,top=19mm,bottom=21mm]{geometry}
\usepackage[colorlinks=true,linkcolor=blue!55!black,citecolor=blue!55!black,
            urlcolor=blue!55!black]{hyperref}

\newtheorem{theorem}{Theorem}[section]
\newtheorem{proposition}[theorem]{Proposition}
\newtheorem{lemma}[theorem]{Lemma}
\newtheorem{corollary}[theorem]{Corollary}
\newtheorem{principle}[theorem]{Principle}
\theoremstyle{definition}
\newtheorem{definition}[theorem]{Definition}
\theoremstyle{remark}
\newtheorem{remark}[theorem]{Remark}
\newtheorem{example}[theorem]{Example}

\newcommand{\Z}{\mathbb Z}
\newcommand{\Q}{\mathbb Q}
\newcommand{\U}{\operatorname U}
\newcommand{\aug}{\varpi}
\newcommand{\dlt}{\delta}
\newcommand{\gm}{\gamma}
\newcommand{\Sym}{\operatorname{Sym}}
\newcommand{\Der}{\operatorname{Der}}
\newcommand{\gr}{\operatorname{gr}}
\newcommand{\cl}{\operatorname{cl}}
\newcommand{\ad}{\operatorname{ad}}
\newcommand{\ab}{\mathrm{ab}}

\setlist{itemsep=3pt,topsep=5pt}

\title[Applications of dimension centrality]
{Applications of the adjoint-module identity\\
for Lie dimension subrings}
\author{}
\date{}

\begin{document}

\begin{abstract}
The identity
\[
                 [\dlt_n(L),L]=\gm_{n+1}(L)
\]
holds for every Lie ring, with no metabelian, finiteness, torsion, or PBW
hypothesis.  We develop its strongest immediate consequences.  The more
fundamental relative identity is
\[
 [A,\dlt_n(L)]\subseteq[A,{}_nL]
\]
for every ideal $A$.  It implies, in particular,
\([\dlt_m(L),\dlt_n(L)]\subseteq\gm_{m+n}(L)\): all brackets between
dimension defects are already classical lower-central terms.

For finite Lie rings this gives an atomic reduction: every nonzero dimension
quotient has a quotient witness which is a monolithic finite $p$-Lie ring with
cyclic center.  This substantially shortens the reduction in the existing
odd-dimensional metabelian proof.  Combining that proof with the new identity
also gives, for every finite nilpotent Lie ring $L$ of class at most $c\geq2$,
\[
                  \dlt_{2c-1}(L)\subseteq Z(L)\cap L''.
\]
Thus metabelianity is no longer needed to localize the possible obstruction: it
is needed only to kill the central second-derived remainder.  For finite
$2$-primary Lie rings a counterexample reduces further to one with cyclic
center and an order-two socle lying simultaneously in
\(\dlt_{2c-1}\), the last lower-central term, and $L''$.

We also give a precise audit of the remaining uses of metabelianity.  The
$\Sym(L^{\ab})$-module structure on $L'$ is equivalent to metabelianity, while
$L'/L''$ has this structure for every $L$.  A bracket-tree lemma identifies the
additional nonmetabelian PBW packets that a genuine extension of the proof
would have to control.
\end{abstract}

\maketitle

\section{The adjoint-module mechanism in its relative form}

We first work over an arbitrary commutative ring $k$ with identity.  Let $L$ be
a Lie algebra over $k$, let $\iota:L\to\U_k(L)$ be the canonical map, and let
$\aug=\ker(\U_k(L)\to k)$ be the augmentation ideal.  No injectivity of
$\iota$ and no PBW theorem will be used.  Put
\[
 \dlt_n(L)=\iota^{-1}(\aug^n),\qquad
 \gm_1(L)=L,\qquad \gm_{n+1}(L)=[\gm_n(L),L].
\]
For an ideal $A\lhd L$, use the notation
\[
 [A,{}_0L]=A,\qquad [A,{}_{r+1}L]=[[A,{}_rL],L].
\]

\begin{lemma}[adjoint augmentation identity]\label{lem:relative-module}
Every ideal $A\lhd L$ is a right $\U_k(L)$-module under the adjoint action,
and for every $n\geq1$ one has
\begin{equation}\label{eq:relative-module}
                    A\cdot\aug^n=[A,{}_nL].
\end{equation}
\end{lemma}

\begin{proof}
On the tensor algebra $T_k(L)$ define
\[
 a\cdot(x_1\cdots x_r)=[a,x_1,\ldots,x_r]
\]
with left-normed brackets.  For $a\in A$ and $x,y\in L$, Jacobi gives
\[
 [[a,x],y]-[[a,y],x]-[a,[x,y]]=0.
\]
Consequently the elements $xy-yx-[x,y]$ act trivially, and the action factors
through $\U_k(L)$.  Since $A$ is an ideal, it is stable under this action.

The additive $k$-module $\aug^n$ is generated by the images of monomials
$x_1\cdots x_r$ with $r\geq n$.  Indeed, $\aug$ is additively generated by
nonempty monomials; multiplication gives one inclusion, and a monomial of
length $r\geq n$ can be split into $n$ nonempty factors for the other.  It
follows that $A\cdot\aug^n\subseteq[A,{}_nL]$, because the relative lower
central series is descending.  Conversely, the commutators of exact length
$n$ after the initial entry generate $[A,{}_nL]$ and are obtained by acting
with monomials $x_1\cdots x_n\in\aug^n$.
\end{proof}

\begin{theorem}[relative dimension-centrality]\label{thm:relative}
For every ideal $A\lhd L$ and every $n\geq1$,
\begin{equation}\label{eq:relative}
                   [A,\dlt_n(L)]\subseteq[A,{}_nL].
\end{equation}
In particular, for every $L$ and $n$,
\begin{equation}\label{eq:main-centrality}
                   [\dlt_n(L),L]=\gm_{n+1}(L).
\end{equation}
\end{theorem}

\begin{proof}
If $a\in A$ and $x\in\dlt_n(L)$, then
\[
 [a,x]=a\cdot\iota(x)\in A\cdot\aug^n=[A,{}_nL]
\]
by Lemma~\ref{lem:relative-module}.  This proves \eqref{eq:relative}.
Taking $A=L$ gives
$[L,\dlt_n(L)]\subseteq\gm_{n+1}(L)$.  The standard inclusion
$\gm_n(L)\subseteq\dlt_n(L)$ gives the reverse inclusion
\[
 \gm_{n+1}(L)=[\gm_n(L),L]
              \subseteq[\dlt_n(L),L].
\]
Signs do not matter for the generated submodules.
\end{proof}

\begin{remark}
The proof is base-ring independent and uses only the universal property of
$\U_k(L)$.  Thus the result is stronger than a statement about integral Lie
rings: it holds over every commutative ground ring, including rings of positive
characteristic and rings with torsion.
\end{remark}

\begin{remark}[why the mechanism is Lie-specific]
The decisive point is that the underlying additive group of $L$ carries the
adjoint action by linear endomorphisms, so that the action extends to
$\U(L)$ and one can form $A\cdot\aug^n$.  A group acting on itself by
conjugation is not a module over its integral group ring.  Thus the identity
$A\cdot\aug^n=[A,{}_nL]$, on which the proof rests, has no literal group
counterpart.  This pinpoints the structural source of the difference from the
group dimension filtration.
\end{remark}

\section{Strong centrality and square-zero dimension defects}

The case $A=L$ is only the first consequence of
Theorem~\ref{thm:relative}.

\begin{corollary}[the entire tail agrees]\label{cor:tails}
For $n,r\geq1$,
\begin{equation}\label{eq:tails}
                  [\dlt_n(L),{}_rL]=\gm_{n+r}(L).
\end{equation}
Thus $\dlt_n(L)$ and $\gm_n(L)$ generate exactly the same lower-central
tail after the zeroth stage.
\end{corollary}

\begin{proof}
The case $r=1$ is \eqref{eq:main-centrality}; bracket repeatedly with $L$.
\end{proof}

\begin{corollary}[mixed commutator theorem]\label{cor:mixed}
For all $m,n\geq1$,
\begin{align}
 [\gm_m(L),\dlt_n(L)]&\subseteq\gm_{m+n}(L),
                                                        \label{eq:gamma-delta}\\
 [\dlt_m(L),\dlt_n(L)]&\subseteq\gm_{m+n}(L).
                                                        \label{eq:delta-delta}
\end{align}
More generally, for $s\geq2$ and $n_1,\ldots,n_s\geq1$,
\begin{equation}\label{eq:iterated-delta}
 [\dlt_{n_1}(L),\dlt_{n_2}(L),\ldots,\dlt_{n_s}(L)]
       \subseteq\gm_{n_1+\cdots+n_s}(L).
\end{equation}
\end{corollary}

\begin{proof}
Apply \eqref{eq:relative} with $A=\gm_m(L)$ to obtain
\eqref{eq:gamma-delta}.  Equation \eqref{eq:main-centrality} shows that
$\dlt_m(L)$ is an ideal: its bracket with $L$ is
$\gm_{m+1}(L)\subseteq\gm_m(L)\subseteq\dlt_m(L)$.  Hence
Theorem~\ref{thm:relative} with $A=\dlt_m(L)$ and
Corollary~\ref{cor:tails} give
\[
 [\dlt_m,\dlt_n]\subseteq[\dlt_m,{}_nL]=\gm_{m+n}.
\]
Iteration using \eqref{eq:gamma-delta} proves
\eqref{eq:iterated-delta}.
\end{proof}

The inclusion \eqref{eq:delta-delta} is strictly stronger than the usual
$N$-series inclusion $[\dlt_m,\dlt_n]\subseteq\dlt_{m+n}$: the bracket
lands not merely in the dimension filtration but already in the classical
lower-central filtration.

\begin{corollary}[the defect has zero bracket]\label{cor:square-zero}
The bracket of $L$ induces well-defined maps
\[
 \dlt_m(L)/\gm_m(L)\ \otimes_k\
 \dlt_n(L)/\gm_n(L)
 \longrightarrow \dlt_{m+n}(L)/\gm_{m+n}(L),
\]
and every one of these maps is zero.  In particular, the graded dimension
defect
\[
                 \mathcal D(L)=\bigoplus_{n\geq1}
                 \dlt_n(L)/\gm_n(L)
\]
is an abelian graded Lie algebra with the bracket induced from $L$.
Moreover, $\dlt_n(L)/\gm_n(L)$ is a central ideal of $L/\gm_n(L)$.
\end{corollary}

\begin{proof}
Changing a representative in degree $m$ by an element of $\gm_m$ changes
the bracket by an element of $\gm_{m+n}$ by
\eqref{eq:gamma-delta}; similarly in degree $n$.  The resulting bracket is
zero by \eqref{eq:delta-delta}.  Finally,
$[\dlt_n,L]=\gm_{n+1}\subseteq\gm_n$.
\end{proof}

\begin{proposition}[virtual lower-central degree]\label{prop:derivations}
Every $a\in\dlt_n(L)$ acts on $\gr_{\gm}L$ as a homogeneous derivation of
degree $n$:
\[
 \overline{x}\in\gm_m/\gm_{m+1}
 \longmapsto [a,x]+\gm_{m+n+1}\in\gm_{m+n}/\gm_{m+n+1}.
\]
This action depends only on the class of $a$ modulo $\dlt_{n+1}(L)$ and
therefore gives a natural map
\begin{equation}\label{eq:graded-derivation}
 \dlt_n(L)/\dlt_{n+1}(L)\longrightarrow
 \Der_n(\gr_{\gm}L).
\end{equation}
\end{proposition}

\begin{proof}
Equation \eqref{eq:gamma-delta} gives the asserted degree and makes the map
independent of the representative of $x$.  Replacing $a$ by an element of
$\dlt_{n+1}$ raises the image by one additional lower-central degree.  The
derivation identity is Jacobi.
\end{proof}

Thus a dimension element of level $n$, even if it does not belong to
$\gm_n$, acts on every classical lower-central layer exactly as an element of
weight $n$ would act.

\begin{proposition}[quotient lifting principle]\label{prop:quotient-lifting}
Let $I\lhd L$.  If
\(\dlt_n(L/I)=\gm_n(L/I)\), then
\begin{equation}\label{eq:quotient-lifting}
                    \dlt_n(L)\subseteq\gm_n(L)+I.
\end{equation}
If, more strongly, $\dlt_n(L/I)=0$, then
\begin{equation}\label{eq:zero-quotient-lifting}
 \dlt_n(L)\subseteq I,\qquad
 \gm_{n+1}(L)\subseteq[I,L].
\end{equation}
If in addition $\gm_{n+1}(L)=0$, then
$\dlt_n(L)\subseteq I\cap Z(L)$.
\end{proposition}

\begin{proof}
Functoriality sends $\dlt_n(L)$ into $\dlt_n(L/I)$, while
\[
 \gm_n(L/I)=(\gm_n(L)+I)/I.
\]
This proves \eqref{eq:quotient-lifting} and the first inclusion in
\eqref{eq:zero-quotient-lifting}.  The latter inclusion and
Theorem~\ref{thm:relative} give
\[
 \gm_{n+1}(L)=[\dlt_n(L),L]\subseteq[I,L].
\]
The final assertion follows again from
$[\dlt_n(L),L]=\gm_{n+1}(L)$.
\end{proof}

\subsection{Consequences in nilpotency class \texorpdfstring{$c$}{c}}

Let $Z_r(L)$ denote the upper central series, with $Z_1(L)=Z(L)$.

\begin{corollary}\label{cor:nilpotent-structure}
If $L$ has nilpotency class at most $c$, then:
\begin{enumerate}[label=\textup{(\roman*)}]
\item for $1\leq n\leq c$,
\[
             \gm_n(L)\subseteq\dlt_n(L)\subseteq Z_{c+1-n}(L);
\]
\item $\dlt_n(L)\subseteq Z(L)$ for every $n\geq c$;
\item $[\dlt_m(L),\dlt_n(L)]=0$ whenever $m+n\geq c+1$;
\item for every $j\geq2$, the lower central series of the Lie ring
$\dlt_n(L)$ satisfies
\[
             \gm_j(\dlt_n(L))\subseteq\gm_{jn}(L),
\]
so $\dlt_n(L)$ is abelian whenever $2n\geq c+1$;
\item $\dlt_n(L)$ centralizes $\gm_m(L)$ whenever $m+n\geq c+1$.
\end{enumerate}
\end{corollary}

\begin{proof}
For (i), equation \eqref{eq:tails} gives
$[\dlt_n,{}_{c+1-n}L]=\gm_{c+1}=0$, which is the defining upper-central
condition.  Part (ii) is \eqref{eq:main-centrality}, and (iii)--(v) follow
from Corollary~\ref{cor:mixed}.  The assertion about
$\gm_j(\dlt_n)$ follows by induction from \eqref{eq:iterated-delta}.
\end{proof}

More generally, the image of $\dlt_n(L)$ in $L/\gm_{n+r}(L)$ lies in its
$r$th upper center; for $r\geq2$, equation \eqref{eq:tails} shows that this
bound is exact whenever $\gm_{n+r-1}/\gm_{n+r}$ is nonzero.

\begin{corollary}[dimension residual]\label{cor:residual}
Put
\[
 \dlt_\omega(L)=\bigcap_{n\geq1}\dlt_n(L),\qquad
 \gm_\omega(L)=\bigcap_{n\geq1}\gm_n(L).
\]
Then
\[
                  [\dlt_\omega(L),L]\subseteq\gm_\omega(L).
\]
Consequently, if $L$ is residually nilpotent, then
$\dlt_\omega(L)\subseteq Z(L)$; if it is also centerless, its dimension
filtration is separated: $\dlt_\omega(L)=0$.
\end{corollary}

\begin{proof}
An element of $\dlt_\omega(L)$ belongs to every $\dlt_n(L)$, so its bracket
with $L$ belongs to every $\gm_{n+1}(L)$ by
\eqref{eq:main-centrality}.
\end{proof}

\section{An atomic testing theorem for finite dimension quotients}

From this point onward, $\U(L)$ and $\dlt_n(L)$ refer to the integral
enveloping algebra of a Lie ring over $\Z$.

The centrality of the defect has a strong reduction-theoretic consequence.
Recall that a nonzero Lie ring is \emph{monolithic} if it has a unique minimal
nonzero ideal, equivalently if some nonzero ideal is contained in every
nonzero ideal.
Here a finite $p$-Lie ring means a Lie ring whose additive group is a finite
$p$-group; it need not be an algebra over $\mathbb F_p$.

We use two elementary facts.  First, dimension subrings are functorial: a Lie
homomorphism $f:L\to M$ satisfies
\begin{equation}\label{eq:functoriality}
                       f(\dlt_n(L))\subseteq\dlt_n(M),
\end{equation}
because $\U(f)$ preserves augmentation powers.  Second, if $B$ is an abelian
Lie ring, then $\U(B)=\Sym(B)$ is graded and
\begin{equation}\label{eq:abelian-dimension}
                       \dlt_n(B)=0\qquad(n\geq2).
\end{equation}

\begin{theorem}[atomic quotient theorem]\label{thm:atomic}
Let $L$ be a finite Lie ring and let $n\geq2$.  If
\[
                         \dlt_n(L)/\gm_n(L)\ne0,
\]
then there are a prime $p$, a quotient $Q$ of $L$, and an element
$0\ne a\in Q$ of order $p$ with all of the following properties:
\begin{enumerate}[label=\textup{(\roman*)}]
\item $\gm_n(Q)=0$ and $0\ne a\in\dlt_n(Q)$;
\item $Q$ is a nonabelian finite nilpotent $p$-Lie ring of class at most
$n-1$;
\item $A=\langle a\rangle$ is contained in every nonzero ideal of $Q$;
in particular, $Q$ is monolithic with monolith $A$;
\item $Z(Q)$ is a cyclic $p$-group;
\item if $d=\cl(Q)$, then $A\subseteq\gm_d(Q)$ and
$\gm_d(Q)$ is a cyclic $p$-group;
\item $\dlt_n(Q)$ is a nonzero cyclic subgroup of $Z(Q)$.
\end{enumerate}
\end{theorem}

\begin{proof}
Choose an element of prime order $p$ in the nonzero finite abelian group
$\dlt_n(L)/\gm_n(L)$.  In
\(M=L/\gm_n(L)\), represent it by
$0\ne a\in\dlt_n(M)$ of order $p$.  Here $\gm_n(M)=0$, so
Theorem~\ref{thm:relative} gives
\[
                         [a,M]=0.
\]
Choose, among the ideals of the finite ring $M$ which do not contain $a$, a
maximal one $K$, and put $Q=M/K$.  We continue to write $a$ for its nonzero
image.  Functoriality preserves membership in $\dlt_n$.  If $J$ is a nonzero
ideal of $Q$ not containing $a$, its inverse image would be an ideal properly
containing $K$ and still avoiding $a$, a contradiction.  This proves (i) and
(iii).  The quotient $Q$ is nilpotent of class at most $n-1$.  It cannot be
abelian, by \eqref{eq:abelian-dimension}.

For each prime $q$, the $q$-primary component $Q_q$ of the additive group of
$Q$ is an ideal: if $q^e x=0$, then $q^e[x,y]=0$ for every $y$.  If
$q\ne p$ and $Q_q\ne0$, this ideal cannot contain the element $a$ of order
$p$, contradicting (iii).  Hence $Q=Q_p$.

Every additive subgroup of $Z(Q)$ is an ideal, so every nonzero such subgroup
contains $A$.  Thus the finite abelian $p$-group $Z(Q)$ has a unique subgroup
of order $p$.  A finite abelian $p$-group with this property is cyclic.  This
proves (iv).  The last nonzero lower-central term $\gm_d(Q)$ is a nonzero
central ideal, so it contains $A$ and is cyclic.  Finally,
$\dlt_n(Q)$ is nonzero and central, and is therefore a subgroup of the cyclic
group $Z(Q)$.
\end{proof}

\begin{corollary}[finite testing principle]\label{cor:testing}
For a fixed $n\geq2$, the following are equivalent:
\begin{enumerate}[label=\textup{(\roman*)}]
\item $\dlt_n(L)=\gm_n(L)$ for every finite Lie ring $L$;
\item $\dlt_n(Q)=0$ for every monolithic finite $p$-Lie ring $Q$, for every
prime $p$, such that $\gm_n(Q)=0$ and $Z(Q)$ is cyclic.
\end{enumerate}
It is enough in (ii) to test rings of class at most $n-1$ whose last
lower-central term contains the order-$p$ monolith.
\end{corollary}

\begin{proof}
Only (ii)$\Rightarrow$(i) needs proof, and its contrapositive is
Theorem~\ref{thm:atomic}.
\end{proof}

This is a genuine gain over merely saying that the quotient
$\dlt_n/\gm_n$ is central.  It replaces an arbitrary finite counterexample by
one which is simultaneously prime-primary, nilpotent, monolithic, and has
cyclic center.  In particular, a last-layer cyclicity assumption can often be
obtained from the new theorem rather than imposed at the outset.

\section{A shorter architecture for the odd metabelian proof}

Continue with integral Lie rings.  Write $c=n+1$, so the index in the existing
odd-dimensional proof is
\[
                         2n+1=2c-1.
\]
The long PBW calculation in
\texttt{metabelian\_dimension\_vanishing.tex} proves, for every $d\geq2$,
the following cyclic-top implication:
\begin{equation*}\tag{$\mathrm{CT}_d$}
\begin{gathered}
 M\text{ finite metabelian of class }d,\quad
 \gm_d(M)\text{ a cyclic }2\text{-group},\\
 x\in\gm_d(M)\cap\dlt_{2d-1}(M)
 \quad\Longrightarrow\quad x=0.
\end{gathered}
\end{equation*}
This is the ``top-layer implication'' proved by the character and PBW assembly.
Passi--Sicking's
other theorem \cite{PS} gives
\begin{equation}\label{eq:PS}
                  2\dlt_r(M)\subseteq\gm_r(M)
                  \qquad(M\text{ metabelian}).
\end{equation}

\begin{proposition}[global vanishing from the cyclic core]
\label{prop:short-metabelian}
Assume $(\mathrm{CT}_d)$ for every $2\leq d\leq c$.  Then every finite
metabelian Lie ring $L$ of class at most $c$ satisfies
\begin{equation}\label{eq:global-metabelian}
                         \dlt_{2c-1}(L)=0.
\end{equation}
\end{proposition}

\begin{proof}
Suppose $0\ne a\in\dlt_{2c-1}(L)$.  Since
$\gm_{2c-1}(L)=0$, \eqref{eq:PS} gives $2a=0$.  Also
Theorem~\ref{thm:relative} gives
\[
 [a,L]\subseteq[\dlt_{2c-1}(L),L]=\gm_{2c}(L)=0.
\]
Apply the maximal-ideal construction in the proof of
Theorem~\ref{thm:atomic}, preserving $a$.  It produces a metabelian quotient
$Q$ which is a finite $2$-Lie ring with cyclic center.  Let
$d=\cl(Q)$.  The ring $Q$ is not abelian, and its last term
$\gm_d(Q)$ is cyclic and contains $a$.  Functoriality and the descending
dimension filtration give
\[
        a\in\dlt_{2c-1}(Q)\subseteq\dlt_{2d-1}(Q),
\]
because $d\leq c$.  This contradicts $(\mathrm{CT}_d)$.
\end{proof}

\begin{remark}[what changed in the existing proof]
Proposition~\ref{prop:short-metabelian} replaces the old induction through
lower nilpotency classes and the separate cyclic-last-layer reduction.  A
hypothetical obstruction is central from the beginning; a maximal quotient
which preserves an order-two obstruction is automatically monolithic, and
monolithicity forces its center and last lower-central term to be cyclic
$2$-groups.  The PBW calculation itself is unchanged.

The Passi--Sicking inclusion \eqref{eq:PS} is still used, but now for exactly
one purpose: it forces the prime of the atomic quotient to be $2$.  Their
previously metabelian centrality assertion is completely replaced by
Theorem~\ref{thm:relative}.
\end{remark}

\subsection{What survives without metabelianity}

The global metabelian theorem has a strong consequence for every finite
nilpotent Lie ring.

\begin{theorem}[central second-derived remainder]\label{thm:central-remainder}
Let $L$ be a finite Lie ring of nilpotency class at most $c\geq2$.  Then
\begin{equation}\label{eq:central-remainder}
                    \dlt_{2c-1}(L)\subseteq Z(L)\cap L''.
\end{equation}
No metabelianity or primary hypothesis on $L$ is required.
\end{theorem}

\begin{proof}
Let $\overline L=L/L''$.  It is a finite metabelian Lie ring of class at most
$c$, so Proposition~\ref{prop:short-metabelian} gives
$\dlt_{2c-1}(\overline L)=0$.  Apply
Proposition~\ref{prop:quotient-lifting} with $I=L''$.  The class bound gives
\[
                         \gm_{2c}(L)=0,
\]
so the final assertion of that proposition is precisely
\eqref{eq:central-remainder}.
\end{proof}

The same argument has a useful variable-index form.  If the metabelianization
$L/L''$ has class $d\geq2$, then
\begin{equation}\label{eq:variable-remainder}
 \dlt_{2d-1}(L)\subseteq L'',\qquad
 \dlt_{\max\{c,2d-1\}}(L)\subseteq Z(L)\cap L''.
\end{equation}
If $L/L''$ is abelian, use $\dlt_2(L/L'')=0$ in place of the first index.

Theorem~\ref{thm:central-remainder} is the sharpest unconditional weakening of
metabelianity supplied by the two proofs alone.  It says that the existing
metabelian machinery already kills the entire image of a putative obstruction
in $L/L''$, while the new theorem makes the remaining error central.  What is
not yet proved is that the central group $Z(L)\cap L''$ contains no such
dimension element.

\begin{corollary}[the exact $2$-primary target]\label{cor:2-primary-target}
Suppose that the conclusion $\dlt_{2c-1}=0$ fails for some finite
$2$-primary Lie ring of class at most $c$.  Then it fails for a finite
$2$-primary quotient $Q$ of some class $2\leq d\leq c$ for which:
\begin{enumerate}[label=\textup{(\roman*)}]
\item $Q$ is monolithic and $Z(Q)$ is a cyclic $2$-group;
\item its monolith $A$ has order $2$;
\item
\begin{equation}\label{eq:atomic-2-target}
 0\ne A\subseteq
 \dlt_{2d-1}(Q)\cap\gm_d(Q)\cap Q''\cap Z(Q).
\end{equation}
\end{enumerate}
Consequently, proving that \eqref{eq:atomic-2-target} is impossible is
sufficient to prove $\dlt_{2c-1}=0$ for all finite $2$-primary Lie rings.
\end{corollary}

\begin{proof}
Choose an element of order $2$ in the nonzero central group
$\dlt_{2c-1}(L)$ and apply the atomic construction.  The resulting quotient
is still $2$-primary and has cyclic center; its monolith $A$ lies in its last
lower-central term.  Since
$\dlt_{2c-1}(Q)\subseteq\dlt_{2d-1}(Q)$, it lies in the displayed dimension
subring.  Apply Theorem~\ref{thm:central-remainder} to $Q$ with its actual
class $d$ to obtain $A\subseteq Q''$.
\end{proof}

This makes the role of $2$-primarity precise.  It can replace
\eqref{eq:PS} in the \emph{reduction}: an order-two obstruction and a cyclic
$2$-primary last layer are then automatic after passage to an atomic quotient.
It does not, by itself, supply the missing proof inside that quotient.

\begin{principle}[$2$-primary globalization]\label{prin:2-globalization}
Let $\mathcal C$ be a class of finite $2$-primary nilpotent Lie rings closed
under quotients.  Suppose that for every $Q\in\mathcal C$ of exact class
$d\geq2$ whose last lower-central term is cyclic, one has the top-layer
implication
\begin{equation}\label{eq:general-CT}
        \gm_d(Q)\cap\dlt_{2d-1}(Q)=0.
\end{equation}
Then every $L\in\mathcal C$ of class at most $c$ satisfies
\[
                         \dlt_{2c-1}(L)=0.
\]
No metabelian hypothesis and no Passi--Sicking exponent theorem is used in
this passage from the cyclic-top case to the general case.
\end{principle}

\begin{proof}
If $\dlt_{2c-1}(L)\ne0$, choose an order-two element in it.  It is central by
\eqref{eq:main-centrality}.  The atomic construction gives a quotient
$Q\in\mathcal C$ of some exact class $d\leq c$ whose last lower-central term
is cyclic and contains the surviving element.  That element belongs to
$\dlt_{2c-1}(Q)\subseteq\dlt_{2d-1}(Q)$, contradicting
\eqref{eq:general-CT}.
\end{proof}

\subsection{A module-separation criterion for the remaining center}

The adjoint module detects all commutators but necessarily annihilates the
center.  The following observation gives an exact representation-theoretic
way to finish an atomic case.

\begin{proposition}[module separation]\label{prop:module-separation}
Let $C$ be an additive subgroup of $L$, let $r\geq1$, and let
$V$ be a right $\U(L)$-module such that
\begin{equation}\label{eq:loewy-bound}
                          V\cdot\aug(L)^r=0.
\end{equation}
If the action map
\[
 C\longrightarrow\operatorname{End}_{\Z}(V),
 \qquad a\longmapsto(v\mapsto v\cdot\iota(a)),
\]
is injective, then
\begin{equation}\label{eq:module-separation}
                           C\cap\dlt_r(L)=0.
\end{equation}
More generally, a family of modules satisfying \eqref{eq:loewy-bound} and
jointly separating the elements of $C$ has the same conclusion.
\end{proposition}

\begin{proof}
If $a\in\dlt_r(L)$, then $\iota(a)\in\aug(L)^r$, so
$v\cdot\iota(a)=0$ for every $v\in V$ by \eqref{eq:loewy-bound}.  Injectivity
on $C$ forces $a=0$.
\end{proof}

For a class-$c$ Lie ring, the adjoint module $L$ satisfies
$L\cdot\aug^c=\gm_{c+1}=0$, but its action map kills $Z(L)$.  Thus it proves
centrality and then stops for an unavoidable reason.  In the atomic situation
of Corollary~\ref{cor:2-primary-target}, it would be enough to construct one
module $V$ with
\[
 V\cdot\aug^{,2d-1}=0
 \quad\text{and}\quad
 V\cdot A\ne0.
\]
Because $A$ has order two and is the unique minimal ideal, this is a sharply
posed module-construction problem.  It is an alternative to adding all the
nonmetabelian PBW packets described below.

\section{Exactly where metabelianity remains}

The centrality theorem comes from treating $L$ as a $\U(L)$-module.  This does
not automatically turn $L'$ into a module over
$\Sym(L^{\ab})$.  The distinction is exact, not cosmetic.

\begin{proposition}[the symmetric-module obstruction]\label{prop:sym-iff}
Let $L$ be a Lie ring and give $L'$ its adjoint right $\U(L)$-module
structure.
\begin{enumerate}[label=\textup{(\roman*)}]
\item The action on $L'$ factors through
\[
             \U(L^{\ab})=\Sym(L^{\ab})
\]
if and only if $L''=0$.
\item For every $L$, the quotient $L'/L''$ is canonically a
$\Sym(L^{\ab})$-module.
\item The failure of the operators defined by $x,y\in L$ to commute on
$m\in L'$ is exactly
\begin{equation}\label{eq:curvature}
 (m\cdot x)\cdot y-(m\cdot y)\cdot x
                     =m\cdot[x,y]=[m,[x,y]]\in L''.
\end{equation}
\end{enumerate}
\end{proposition}

\begin{proof}
The action factors through $L/L'$ precisely when every element of $L'$ acts
trivially on $L'$, which says $[L',L']=L''=0$.  This proves (i).  Modulo
$L''$, elements of $L'$ act trivially and \eqref{eq:curvature} vanishes, which
proves (ii).  Formula \eqref{eq:curvature} is the defining enveloping-algebra
relation, or directly Jacobi.
\end{proof}

Thus the old symmetric-power characters always exist on the
metabelianization, and this is precisely why they prove
Theorem~\ref{thm:central-remainder}.  To go farther, one must control the
$L''$-valued curvature \eqref{eq:curvature}; centrality of the ultimate
dimension element does not make that curvature vanish.

\begin{example}[$2$-primary does not imply the symmetric module condition]
\label{ex:ut5}
Let $R=\Z/2^e$ and let $L=\mathfrak n_5(R)$ be the Lie ring of strictly upper
triangular $5\times5$ matrices, with bracket $[x,y]=xy-yx$.  It is a finite
$2$-primary Lie ring of class $4$.  With $E_{ij}$ denoting a matrix unit,
\[
 [E_{12},E_{23}]=E_{13},\qquad
 [E_{34},E_{45}]=E_{35},\qquad
 [E_{13},E_{35}]=E_{15}\ne0.
\]
Hence $L''\ne0$.  More explicitly, for
$m=E_{35}\in L'$, $x=E_{12}$ and $y=E_{23}$, the operator commutator in
\eqref{eq:curvature} is
\[
                         [m,[x,y]]=-E_{15}\ne0.
\]
Thus neither $2$-primarity nor even the fact that this particular second
derived term is central makes $L'$ a $\Sym(L^{\ab})$-module.  This example is
not asserted to be a dimension-quotient counterexample; it isolates the
failure of the proposed module shortcut.
\end{example}

In the existing proof, metabelianity has the following remaining concrete
uses.
\begin{enumerate}
\item The relatively free model has $M=F'$ with $[M,M]=0$, so $M$ is an
honest $\Sym(F^{\ab})$-module.  Proposition~\ref{prop:sym-iff} shows that this
exact statement cannot hold in a genuinely nonmetabelian ring.
\item The maps
\[
 \Theta_k:\gm_k/\gm_{k+1}\otimes
       \Sym^{m_k-1}(L/\gm_k)\longrightarrow\gm_c
\]
are well defined and symmetric because every Jacobi interchange error has
two derived inputs.  Without metabelianity those errors land in $L''$, not
in zero.
\item In the quadratic PBW block, interchanging two derived factors creates
a primitive commutator in $L''$.  Those were the ``silent vertical
corrections'' in the metabelian proof and are no longer silent.
\item In the final PBW assembly, the assertion that a nonzero comb has at
most one $M$-input deletes all bracket trees with two derived branches.  The
next lemma identifies their values exactly.
\end{enumerate}
The new theorem removes none of these four identities.  It instead proves
that, after all metabelian packets have been killed, their total possible
error is confined to the central group $Z(L)\cap L''$.

\section{The missing nonmetabelian packets: a combinatorial localization}

The extra terms are not arbitrary.  They are indexed by branching in a
bracket tree.

\begin{lemma}[derived-branching depth]\label{lem:branching}
Let $T$ be a binary bracket tree evaluated in a Lie ring $L$.  Suppose that
two disjoint rooted subtrees of $T$ evaluate in $L'$.  Let $v$ be their least
common ancestor and let $h$ be the number of edges from $v$ to the root of
$T$.  Then
\begin{equation}\label{eq:branch-depth}
                 T(L)\subseteq[L'',{}_hL].
\end{equation}
In particular, if $L''\subseteq Z(L)$, every such tree with $h\geq1$
vanishes; a nonzero tree with two derived branches must join those branches
at its root.
\end{lemma}

\begin{proof}
Each of the two child branches below $v$ contains one of the distinguished
derived subtrees.  Since $L'$ is an ideal, the value of each complete child
branch lies in $L'$.  Their bracket at $v$ therefore lies in
$[L',L']=L''$.  At every one of the $h$ ancestors, this value is bracketed
with another element of $L$, which proves \eqref{eq:branch-depth}.  If $L''$
is central, the first such further bracket is zero.
\end{proof}

For a ring of class $c$, the root-level errors in the centrally
second-derived case are therefore measured by the finite family of
well-defined pairings
\begin{equation}\label{eq:top-pairings}
 \lambda_i:
 \frac{\gm_i(L)}{\gm_{i+1}(L)}\otimes
 \frac{\gm_{c-i}(L)}{\gm_{c-i+1}(L)}
 \longrightarrow \gm_c(L),
 \qquad \lambda_i(\bar u,\bar v)=[u,v],
\end{equation}
where $2\leq i\leq c-2$.  These maps are well defined because an altered
representative produces an element of $\gm_{c+1}=0$, and their images lie in
$L''\cap\gm_c$.  In a metabelian ring every $\lambda_i$ is zero.  In
Example~\ref{ex:ut5}, the class-four pairing $\lambda_2$ is nonzero.

\begin{proposition}[precise shape of the next calculation]
\label{prop:next-calculation}
Suppose $L$ has class $c$ and $L''\subseteq Z(L)$.  After primitive PBW
terms of total bracket weight $c$ are expanded into binary bracket trees,
every term omitted by the metabelian ``at most one derived input'' rule is
either zero or is a root value represented by one of the pairings
\eqref{eq:top-pairings}.  There is no spectator bracket \emph{above} the
joining of the two derived branches in a nonzero omitted term.
\end{proposition}

\begin{proof}
An omitted tree has two disjoint derived branches.  By
Lemma~\ref{lem:branching}, it vanishes unless their least common ancestor is
the root.  In the latter case the two root branches have weights $i$ and
$c-i$, both at least two, and their value is exactly \eqref{eq:top-pairings}.
An additional spectator would put the least common ancestor below the root
and force zero.
\end{proof}

For general $L$, Lemma~\ref{lem:branching} gives a filtration of the missing
terms by
\[
 J_0=L'',\qquad J_{h+1}=[J_h,L].
\]
A branching vertex creates a value in $J_0$, and its distance from the root
records the $J_h$ in which the final error lies.  This supplies a concrete
extension strategy: the old packets handle trees with at most one derived
branch; new packets must handle the pairings \eqref{eq:top-pairings} and their
$L$-action descendants.  In the atomic $2$-primary reduction, all terminal
values may be evaluated in one cyclic central $2$-group.  What remains open
is the chain-level cancellation of these new packets.

\section{Conclusions and the exact boundary of the result}

The following statements are now proved.
\begin{enumerate}
\item For every Lie algebra over every commutative base ring,
\([\dlt_n,L]=\gm_{n+1}\), and more generally
$[A,\dlt_n]\subseteq[A,{}_nL]$ for every ideal $A$.
\item Dimension elements have genuine lower-central weight under brackets:
$[\dlt_m,\dlt_n]\subseteq\gm_{m+n}$ and
$[\gm_m,\dlt_n]\subseteq\gm_{m+n}$.
\item Every finite nonzero dimension quotient has an atomic witness which is
a monolithic finite $p$-Lie ring with cyclic center.
\item The preliminary induction in the odd metabelian proof can be replaced
by the one-step atomic quotient argument of
Proposition~\ref{prop:short-metabelian}.
\item For every finite nilpotent Lie ring of class at most $c$,
$\dlt_{2c-1}$ is contained in the central second-derived group
$Z(L)\cap L''$.
\item For finite $2$-primary rings, it is enough to exclude the sharply
specified atomic configuration \eqref{eq:atomic-2-target}.
\item The remaining cyclic central socle can be killed by constructing a
$\U(L)$-module of augmentation Loewy length at most $2c-1$ which detects that
socle, as in Proposition~\ref{prop:module-separation}.
\end{enumerate}

What is \emph{not} yet proved is
\(\dlt_{2c-1}(L)=0\) for arbitrary finite $2$-primary, nonmetabelian $L$.
The hypothesis ``$2$-primary'' completely replaces metabelianity in the
reduction to a cyclic central target, but it does not make $L'$ a symmetric
module and it does not annihilate the pairings \eqref{eq:top-pairings}.  The
remaining problem is therefore no longer vague: it is to construct and kill
the $L''$-valued branching packets described in
Lemma~\ref{lem:branching}, ultimately in the cyclic order-two monolith of an
atomic quotient.

\begin{thebibliography}{9}

\bibitem{BP}
L.~Bartholdi and I.~B.~S.~Passi,
\emph{Lie dimension subrings},
Internat. J. Algebra Comput. \textbf{25} (2015), 1301--1325.

\bibitem{PS}
I.~B.~S.~Passi and T.~Sicking,
\emph{Dimension quotients of metabelian Lie rings},
Internat. J. Algebra Comput. \textbf{27} (2017), 251--258;
\href{https://arxiv.org/abs/1602.04919}{arXiv:1602.04919}.

\end{thebibliography}

\end{document}
